\chapter{Introduction and Problem Definition}\label{ch:intro}

\section{Project Background}
Every motor vehicle on an Indian road is tied to a Regional Transport Office. The RTO
is the state-level office that registers vehicles, issues learner and driving
licences, grants permits for commercial transport, collects road tax and fees, and
enforces the Motor Vehicles Act, 1988 by issuing challans (traffic fines). The Ministry
of Road Transport and Highways (MoRTH) lays down the national framework, and the
National Informatics Centre (NIC) builds the common systems, Vahan for vehicle
registration and Sarathi for driving licences, whose data is compiled into state and
national registers \cite{parivahan-about}. What happens at the counter of one office,
however, is still a long chain of small decisions made by different people on
different days: a clerk checks documents, an officer approves, an accountant records
a payment, an inspector passes a test.

Each of those decisions leaves a record, and the records are tightly connected. A
vehicle has an owner, the owner has a licence, the licence has classes, the vehicle
has insurance and a fitness certificate, a violation is committed with a vehicle and
produces a challan, the challan is paid through a payment that can be retried or
refunded. If any link is wrong, for example two current owners for one vehicle or a
challan paid twice, the damage reaches the citizen, the office and the revenue.

This project, the \textbf{RTO Management System}, models that world end to end. It
consists of three parts that depend on each other in a clear order: a MySQL 8 database
of 63 tables that carries the rules, a Spring Boot REST backend that adds security,
workflow and transactions, and a Spring Boot web portal that lets each kind of staff
member and citizen use the system without knowing any SQL.

\section{Problem Statement}
Office work of this kind is usually described as ``data entry'', but it is really a
correctness problem. The system must guarantee, whoever the user is and however many
users act at once, that:
\begin{itemize}
\item a vehicle has exactly one current owner at any moment, and its ownership history is never lost;
\item an appointment slot is never booked beyond its capacity;
\item a fee or a fine is never recorded as paid twice, even if the payment request is repeated;
\item every status change follows a legal path and leaves a trace of who changed it, when and why;
\item each person can see and do only what their role allows, and a citizen sees only their own records.
\end{itemize}
Systems built around paper files, spreadsheets or loosely validated forms cannot make
these guarantees. The problem addressed here is therefore to design and build a
centralised RTO management system whose \emph{database enforces these rules itself},
so that a mistake in an application, a script or a person cannot put wrong data in,
and to demonstrate with tests that the rules actually hold.

\section{Motivation}
We chose an RTO for three reasons. First, the domain is rich enough to need a
real design: the final schema has 63 tables and 91 foreign keys, far beyond the
handful of tables of a typical course exercise, so normalisation, keys, constraints,
indexes and views have to be used for real reasons and not for show. Second, almost
every topic of our database course has a natural home in it. History tables record
status changes, a trigger protects active vehicles from deletion, a stored procedure
transfers ownership in one transaction, row locks protect the last appointment seat,
and window functions rank offices by workload. Third, the problem is familiar. Anyone
who has queued at a transport office, lost a document or paid an agent to push a
form through understands why traceable, rule-driven processing matters.

A further motivation was to test our own work honestly. We wanted claims such as ``two
users cannot book the last seat'' to be backed by an automated test that fires
concurrent requests at a real MySQL server, not merely stated in a report.

\section{Objectives}
The objectives of the project are:
\begin{enumerate}
\item To design a normalised relational schema for the complete RTO domain and to keep the business rules inside the database as constraints, generated columns, triggers and procedures.
\item To build a secure Java REST backend with permissions stored in the database, transactional workflows, safe concurrent updates and an audit trail.
\item To build a web portal in which every role sees only the screens and actions its permissions allow.
\item To show advanced database features working together: indexes, views, triggers, stored procedures, transactions, locking and analytical SQL.
\item To verify the system with automated tests, a crawl of the whole user interface, a load test and measured index experiments, and to report defects found along the way.
\item To apply the concepts of five courses (Object Oriented Programming, Database Management, Software Engineering, Data Structures, and Design and Analysis of Algorithms) to one working system.
\end{enumerate}

\section{Scope of the Project}
\subsection{In scope}
The system covers: citizens, addresses and applicants; offices, counters, departments,
designations and employees with their postings; users, roles and permissions; service
applications with documents, status history and appointments; learner licences,
driving tests, driving licences and licence classes; driving schools and instructors;
vehicles, ownership history and ownership transfers; inspections, fitness, pollution
and insurance certificates and road-tax records; routes and commercial permits;
violations and challans; fee structures, payments, payment attempts and refunds;
complaints, appeals, notifications and an audit log; and a management dashboard.

\subsection{Out of scope}
The project is an academic prototype. It does not connect to Aadhaar, DigiLocker, the
national Vahan and Sarathi registers or a live payment gateway; notifications go to a
mock provider and no real SMS or e-mail is sent; there is no mobile application; and
no claim is made about capacity beyond the single-machine load tests reported in
Chapter~\ref{ch:perf}.

\subsection{Users and environment}
There are seven seeded roles: administrator, RTO officer, counter clerk, inspector,
enforcement officer, accountant and citizen. The backend runs on Java 21 and listens on
port 8000; the portal runs on its own port 8080; both use a MySQL 8.0 server. A normal
web browser is all a user needs.

\section{Proposed Solution}
The solution is organised as three tiers, shown in Chapter~\ref{ch:design}, that are
built and tested in dependency order.

\begin{enumerate}
\item \textbf{Database tier.} Thirteen module scripts (files 02 to 14) create the 63
tables in a fixed dependency order. A second set of scripts adds CHECK constraints, indexes, views, triggers and
stored procedures. A seed script loads realistic data so that every status and filter
can be demonstrated, and a SQL test script checks the enhancements.
\item \textbf{Application tier.} A layered Spring Boot backend: controllers receive
requests; an interceptor enforces authentication and permissions before any controller
code runs; services hold the business rules inside transactions; a small data-access
helper provides locking, paging and auditing.
\item \textbf{Presentation tier.} A server-rendered Thymeleaf portal styled with
Tailwind CSS. It holds the login token in a server-side session, calls the backend,
and shows each user only the menu items their permissions allow.
\end{enumerate}

\section{Novelty and Originality}
Individually, none of the techniques here is new: foreign keys, row locking and JWT
are standard. The contribution lies in how they are combined and in how carefully the
result is checked.
\begin{itemize}
\item \textbf{Two layers guard each important rule.} For instance, the one-current-owner rule is enforced by a unique key on a generated column in the database and, independently, by a locked, ordered check in the service. Either alone would be enough on a good day; together they are enough on a bad one.
\item \textbf{Permissions are data.} Roles and 62 permissions live in tables, so an administrator can change access from the portal without a deployment.
\item \textbf{A deliberate division of labour between triggers and Java.} Triggers were added only where the backend does not already write the row, which avoids duplicate history entries (Section~\ref{sec:trigdesign}).
\item \textbf{Evidence over assertion.} The report states real results, including defects the testing uncovered in the supplied scripts.
\end{itemize}

\section{Feasibility Analysis}
\subsection{Technical feasibility}
Everything used is mature and open source: MySQL 8, Java 21, Spring Boot 3.3.5,
Hibernate, Thymeleaf and Tailwind CSS. The work divides cleanly along the schema and the
REST contract, so the three tiers could be developed in parallel by different team
members and joined later.

\subsection{Economic feasibility}
No licence fees are involved. The whole system runs on one ordinary laptop (Intel Core
i5-1335U, 16~GB RAM, Windows~11) and a free MySQL Community Server. A real deployment
would add hosting, backups and an SMS or e-mail gateway.

\subsection{Operational feasibility}
Staff use a browser. The menu adapts to each role, and error pages are written in plain
language, so a clerk is not shown a stack trace. Reference data such as fees, offices
and slots can be changed from the portal.

\subsection{Schedule feasibility}
The repository history shows the order in which the work was done: the schema first,
then the backend, a first version of the portal, the database enhancements, and finally
the complete portal and the final validation (Section~\ref{sec:timeline}).

\section{Social Relevance and Practical Applications}
Delays, repeated visits and opaque fees at transport offices are everyday experiences.
A traceable system, in which every fee comes from a published table and every status
change has a name and a time attached, lowers the room for informal charges and makes
complaints and appeals easier to resolve. The design could serve as the core of a
state-level pilot, a teaching reference for relational design, or the starting point for
a citizen mobile application, because the REST API is documented with Swagger and is
independent of the portal. It is not a replacement for the national systems and makes
no such claim.

\section{Success Criteria and Key Performance Indicators}
Table~\ref{tab:kpi} lists the measurable criteria we set, how each is measured and
where the result is reported. The results themselves appear in Chapters~\ref{ch:perf}
and~\ref{ch:test}.

\begin{table}[H]
\centering\small
\caption{Key performance indicators and how they are measured}\label{tab:kpi}
\begin{tabularx}{\linewidth}{@{}p{3.1cm}Xp{3.3cm}@{}}
\toprule
\textbf{KPI} & \textbf{Measurement method} & \textbf{Reported in}\\ \midrule
Functional correctness & Share of backend integration tests that pass on a real MySQL server & Section~\ref{sec:unit}\\
Rule enforcement under concurrency & Parallel-request tests for slot booking, payments, refunds, status changes and ownership approval & Section~\ref{sec:conc}\\
Schema integrity & Number of foreign keys with orphan rows; schema audit test; SQL test suite & Sections~\ref{sec:schemaaudit}, \ref{sec:sqltests}\\
UI robustness & Crawl of all reachable pages and edge-case URLs; count of server errors (5xx) & Section~\ref{sec:crawl}\\
Response time & Average and 95th-percentile latency of read operations at 1, 5, 10 and 20 concurrent users & Section~\ref{sec:loadres}\\
Index effectiveness & Lookup time with and without the project's indexes on 10\,000 to 500\,000 rows & Section~\ref{sec:indexres}\\
Access control & Permission checks per role; a low-privilege account cannot reach administrator pages & Sections~\ref{sec:rbactest}, \ref{sec:crawl}\\
\bottomrule
\end{tabularx}
\end{table}
